\documentclass[12pt,a4paper,pl,fleqn]{article}


% \usepackage[english,polish]{babel}                   % ostatni jezyk przewaza w opisach rozdzialow i literatury
\usepackage[utf8]{inputenc}
\usepackage{graphicx}
\usepackage[toc]{appendix}
% \usepackage[T1]{fontenc}



% \documentclass[a4paper,12pt]{article}
% \usepackage[utf8]{inputenc}
\usepackage[T1]{fontenc}
\usepackage[german,polish,english]{babel}
\usepackage{polski}
\usepackage{lmodern}


\setlength{\textheight}{210mm}
\setlength{\textwidth}{145mm}
\setlength{\parindent}{20.0mm}

\renewcommand{\thesection}{\Roman{section}}

\title{Obtaining Schrödinger Equation from Wave Equation}

% \title{Wyprowadzenie równania Schrödingera z równania falowego}

\author{Paweł Jan Piskorz\footnote{Address: Krakowska 55, 31-066 Kraków, Poland; email: paweljanpiskorz@gmail.com}}
\date{}


\begin{document}
\maketitle




% W styczniu 2026 roku mija sto lat kiedy to austriacki fizyk Erwin Schrödinger opublikował równanie nazwane później jego imieniem.
% Równanie Schrödingera może zostać wyprowadzone z równania falowego, hipotezy de~Broglie oraz z wzoru Plancka na energię kwantu fali.

January 2026 marks one hundred years since the Austrian physicist Erwin Schrödinger has published the equation later named after him.
Schrödinger equation can be derived from wave equation, de~Broglie hypothesis and Planck formula for energy of wave quanta.


% Równanie falowe w jednym wymiarze może zostać zapisane jako
The wave equation in one dimension can be written as
\begin{equation}
\frac{1}{v_{p}^{2}} \frac{\partial^{2} \Psi}{\partial t^{2}} = \frac{\partial^{2} \Psi}{\partial x^{2}}
\end{equation}
% gdzie $v_{p}$ oznacza prędkość fazową fali, $t$ oznacza czas natomiast $x$ oznacza współrzędną miejsca, w ktorym obserwujemy
% falę o wielkości $\Psi(x, t)$\@.
where $v_{p}$ is the phase velocity of the wave, $t$ denotes time and $x$ is the coordinate of the point on the $Ox$ axis, where the wave of
magnitude $\Psi(x, t)$ is being observed.



% Hipoteza de~Broglie może zostać wyprowadzona z wzoru Einsteina
De Broglie hypothesis can be derived from Einstein formula
\begin{equation}
E = mc^{2}
\end{equation}
% na energię $E$ cząstki posiadającej masę $m$ i w którym $c$ oznacza prędkość światła w próżni. Otrzymujemy wtedy, że
for energy $E$ of the mass $m$\@. $c$ denotes the speed of light in vacuum. We receive therefrom
\begin{equation}
E = mc^{2} = m c c = p c = p \lambda/T
\end{equation}
where $p$ denotes particle momentum, $\lambda$ is the wave length, and $T$ is the period of the wave in time domain.
On the other hand from the Planck formula we have
% gdzie $p$ oznacza pęd cząstki, $\lambda$ jest długością fali a $T$ jest okresem fali w domenie czasu.
% Z drugiej strony ze wzoru Plancka otrzymujemy, że
\begin{equation}
E = h \nu = h/T
\end{equation}
where $h$ is Planck constant, $\nu$ wave frequency.
Comparing the two above equations for energy $E$ we receive
% gdzie $h$ jest stałą Plancka a $\nu$ jest częstotliwością fali.
% Porównując dwa powyższe równania na energię $E$ dostajemy wzór
\begin{equation}
p \lambda/T = h/T
\end{equation}
and therefrom the de Broglie hypothesis
% a z niego hipotezę de~Broglie
\begin{equation}
p = h / \lambda
\end{equation}
% which relates the particle momentum with its wave length.
that is the formula relating the particle momentum $p$ with the length $\lambda$ of its wave. Such wave we call the {\em de Broglie wave.}
% czyli wzór łączący pęd $p$ cząstki z długością $\lambda$ jej fali, którą nazywamy {\em falą de Broglie.}
% Teraz wyprowadzimy wzór na energię kinetyczną $E - V$ cząstki, gdzie $E$ jest jej energią całkowitą a $V$ jej energią potencjalną.

Now we derive the formula for the kinetic energy $E - V$ of a particle, where $E$ is its total energy and $V$ is its potential energy.
\begin{equation}
E - V = \frac{1}{2} m v_{p}^{2}
\end{equation}
% Kontynuujemy otrzymując
We continue obtaining
\begin{equation}
2 m (E - V) = m^{2} v_{p}^{2} = p^{2}
\end{equation}
\begin{equation}
p = [2 m (E - V)]^{1/2} = h/\lambda
\end{equation}
\begin{equation}
\lambda = \frac{h}{[2 m (E - V)]^{1/2}}
\end{equation}
\begin{equation}
v_{p} = \lambda/T = \lambda \nu = \frac{h \nu}{[2 m (E - V)]^{1/2}} = \frac{E}{[2 m (E - V)]^{1/2}}
\end{equation}
% Zatem
Therefore
\begin{equation}
\frac{1}{v_{p}^{2}} = \frac{2 m (E - V)}{E^{2}}
\end{equation}

The solution of the wave equation can be presented in the form
% Rozwiązanie równania falowego może zostać przedstawione w postaci
\begin{equation}
\Psi =  e^{i \omega t}  \varphi(x) =   e^{i (E/\hbar) t} \varphi(x)     
\end{equation}
where $E = \hbar \omega$ whereby $\hbar = h/2\pi$ and $\omega = 2\pi/T$
\begin{equation}
\frac{\partial \Psi}{\partial t} = i \frac{E}{\hbar} e^{i (E/\hbar) t} \varphi(x)
\end{equation}
\begin{equation}
\frac{\partial^{2} \Psi}{\partial t^{2}} = - \frac{E^{2}}{\hbar^{2}} e^{i (E/\hbar) t} \varphi(x)
\end{equation}
\begin{equation}
\frac{\partial^{2} \Psi}{\partial x^{2}} = e^{i (E/\hbar) t} \, \frac{\partial^{2} \varphi(x)}{\partial x^{2}}
\end{equation}


\begin{equation}
\frac{2m(E - V)}{E^{2}} \frac{\partial^{2} \Psi}{\partial t^{2}} = \frac{\partial^{2} \Psi(x)}{\partial x^{2}}
\end{equation}


\begin{equation}
- \frac{2m(E - V)}{E^{2}} \frac{E^{2}}{\hbar^{2}} e^{i (E/\hbar) t} \varphi(x) = e^{i (E/\hbar) t} \frac{\partial^{2} \varphi(x)}{\partial x^{2}}
\end{equation}

\begin{equation}
- \frac{2m}{\hbar^{2}} (E - V) \varphi(x) = \frac{\partial^{2} \varphi(x)}{\partial x^{2}}
\end{equation}
Therefrom we obtain the time independent Schrödinger equation in one dimension
% Stąd otrzymujemy równanie znane jako równanie Schrödingera w jednym wymiarze niezależne od czasu
\begin{equation}
- \frac{\hbar^{2}}{2 m} \frac{\partial^{2} \varphi(x)}{\partial x^{2}}  + V(x) \varphi(x) = E \varphi(x)
\end{equation}
% które otwiera nowe możliwości opisu fotonów, elektronów, atomów pierwiastków i cząsteczek chemicznych oraz innych układów mikroświata.
which opens up new possibilities for describing photons, electrons, atoms of elements, chemical molecules and other systems of the microworld.






%8888888888888888888888888888888888888888888888888888888888888888

% \begin{thebibliography}{99}
% \bibitem{kn:Borowitz_Book} Borowitz, S. \ 
%                                         (1967) {\em Fundamentals of Quantum Mechanics\/} W.A.~Benjamin, Inc., New York, Amsterdam



% \end{thebibliography}

% \subsection*{}

% \begin{verbatim}
% Pawel J. Piskorz
% \end{verbatim}


\end{document}




